\section{การวิเคราะห์ระบบ}

ผู้ใช้มีเป้าหมายหลักคือการเตรียมสื่อสำหรับเผยแพร่โดยลดการเปิดเผยตัวตนของบุคคล ระบบจึงออกแบบให้เริ่มจากการอัปโหลดไฟล์ แสดงผลการตรวจจับในรูปของรายการบุคคลที่จัดกลุ่มแล้ว เปิดโอกาสให้ผู้ใช้ตัดสินใจว่าจะปกปิดหรือยกเว้นบุคคลใด และแก้ไขบริเวณที่ AI ยังตรวจจับไม่ครอบคลุมก่อนส่งออกไฟล์ ตารางที่ \ref{tab:functional-requirements} แสดงความต้องการเชิงฟังก์ชันของระบบปัจจุบัน

\begin{table}[H]
\centering
\caption{ความต้องการเชิงฟังก์ชันของระบบ}
\label{tab:functional-requirements}
\begin{tabular}{|p{0.13\textwidth}|p{0.42\textwidth}|p{0.30\textwidth}|}
\hline
\textbf{รหัส} & \textbf{ความต้องการ} & \textbf{สถานะในระบบ} \\
\hline
FR-01 & ผู้ใช้อัปโหลดภาพหรือวิดีโอที่รองรับได้ & พัฒนาแล้ว \\
\hline
FR-02 & ระบบสร้างข้อมูลเมตาและเฟรมตัวอย่างของสื่อที่อัปโหลด & พัฒนาแล้ว \\
\hline
FR-03 & ระบบตรวจจับใบหน้าและบันทึกตำแหน่งตามเฟรม & พัฒนาแล้ว \\
\hline
FR-04 & ระบบรวมใบหน้าที่คาดว่าเป็นบุคคลเดียวกันเป็นกลุ่มบุคคล & พัฒนาแล้ว \\
\hline
FR-05 & ผู้ใช้เลือกเบลอทุกคน เฉพาะคน หรือทุกคนยกเว้นคนที่เลือกได้ & พัฒนาแล้ว \\
\hline
FR-06 & ผู้ใช้กำหนด Gaussian Blur, Pixelation หรือ Black Box และปรับความเข้มได้ & พัฒนาแล้ว \\
\hline
FR-07 & ผู้ใช้วาด ย้าย ปรับขนาด ลบ และเลือกติดตามกรอบ Manual Blur ได้ & พัฒนาแล้ว \\
\hline
FR-08 & ระบบสร้างตัวอย่างและส่งออกภาพหรือวิดีโอที่ปกปิดข้อมูลแล้วได้ & พัฒนาแล้ว \\
\hline
FR-09 & ระบบส่งไฟล์ผลลัพธ์ไปยัง Cloudinary ได้เมื่อมีข้อมูลรับรอง & พัฒนาแล้วในลักษณะตัวเลือกเสริม \\
\hline
\end{tabular}
\end{table}

ความต้องการที่ไม่ใช่เชิงฟังก์ชัน ได้แก่ การป้องกันพิกัดกรอบไม่ให้เกินขนาดเฟรม การรองรับงานเบื้องหลังระหว่างตรวจจับวิดีโอ การแสดงข้อความแจ้งสถานะการทำงาน และการคงผลการเลือกปกปิดไว้ในฐานข้อมูล อย่างไรก็ตาม ซอร์สโค้ดปัจจุบันยังต้องปรับปรุงความปลอดภัยสำหรับการใช้งานจริง เช่น การกำหนด CORS ให้จำกัดโดเมน การบังคับตรวจสอบขนาดไฟล์ที่อัปโหลด และการจัดเก็บสถานะงานให้ทนทานต่อการรีสตาร์ตเซิร์ฟเวอร์

\section{การออกแบบข้อมูลและส่วนประกอบ}

ฐานข้อมูลของระบบแบ่งข้อมูลเป็นสามกลุ่ม ได้แก่ ข้อมูลสื่อ ผลการวิเคราะห์ใบหน้า และการตั้งค่าปกปิด ตารางที่ \ref{tab:data-model} สรุป entity สำคัญที่กำหนดด้วย SQLAlchemy

\begin{table}[H]
\centering
\caption{โครงสร้างข้อมูลหลักของระบบ}
\label{tab:data-model}
\begin{tabular}{|p{0.26\textwidth}|p{0.28\textwidth}|p{0.31\textwidth}|}
\hline
\textbf{ตาราง} & \textbf{ข้อมูลสำคัญ} & \textbf{บทบาท} \\
\hline
videos & ชื่อไฟล์, ขนาดเฟรม, FPS, จำนวนเฟรม, สถานะ & เก็บข้อมูลต้นทางของภาพหรือวิดีโอ \\
\hline
preview\_frames & หมายเลขเฟรม เวลา และที่อยู่ของภาพขนาดย่อ & ใช้เลือกและแสดงตัวอย่างเฟรม \\
\hline
detected\_faces & พิกัดกรอบ ค่าความเชื่อมั่น กลุ่มบุคคล และหมายเลขเฟรม & เก็บผลตรวจจับที่ใช้สร้างผลลัพธ์ \\
\hline
identities & ป้ายกำกับ ภาพขนาดย่อ เวกเตอร์ลักษณะเฉลี่ย และจำนวนการปรากฏ & แทนบุคคลแต่ละรายบน Face Map \\
\hline
blur\_selections & รหัสวิดีโอ รหัสกลุ่มบุคคล และสถานะการปกปิด & เก็บการตัดสินใจว่าจะปกปิดบุคคลใด \\
\hline
timeline\_entries & รหัสกลุ่มบุคคล หมายเลขเฟรม และเวลา & แสดงช่วงเวลาที่บุคคลปรากฏ \\
\hline
manual\_blur\_boxes & พิกัด ขนาด เฟรมเริ่มต้น และสถานะการติดตาม & เก็บพื้นที่ปกปิดที่ผู้ใช้กำหนดเอง \\
\hline
\end{tabular}
\end{table}

ส่วนหลัง (backend) แยกความรับผิดชอบเป็นบริการตรวจจับใบหน้า บริการรู้จำใบหน้า บริการติดตามวัตถุ และบริการสร้างผลลัพธ์ การแยกส่วนดังกล่าวช่วยให้เปลี่ยนแบบจำลองหรือปรับวิธีปกปิดได้โดยไม่กระทบ API หลัก ส่วนหน้า (frontend) แบ่งเป็นคอมโพเนนต์สำหรับอัปโหลดสื่อ เลือกบุคคล แสดง Timeline ควบคุมการปกปิดด้วยตนเอง และส่งออกผลลัพธ์

\section{การออกแบบส่วนติดต่อผู้ใช้}

ส่วนติดต่อผู้ใช้จัดลำดับการทำงานให้เริ่มจากการอัปโหลดสื่อ เลือกเครื่องมือ ตรวจสอบผลการตรวจจับ แก้ไขพื้นที่ปกปิด และยืนยันผลก่อนส่งออก ภาพ Demo ในส่วนนี้เรียงตามขั้นตอนการใช้งาน และวางคำอธิบายไว้ใต้ภาพที่เกี่ยวข้องโดยตรง

% ==================================================
% TODO IMAGE: ใส่ภาพ Demo หน้าจอหลักและส่วนอัปโหลดไฟล์
% Caption: หน้าจอหลักและส่วนอัปโหลดไฟล์
% ==================================================
\begin{figure}[H]
\centering
\includegraphics[width=0.94\textwidth]{\detokenize{images-progress/ภาพ demo หน้าแรก (Home).jpg}}
\caption{หน้าจอหลักและส่วนอัปโหลดไฟล์}
\label{fig:demo-home}
\figuredescription{fig:demo-home}{เป็นหน้าจอเริ่มต้นสำหรับเลือกหรือลากไฟล์ภาพและวิดีโอเข้าสู่พื้นที่อัปโหลดของระบบ}
\end{figure}

% ==================================================
% TODO IMAGE: ใส่ภาพ Demo หน้าจอเลือกเครื่องมือและตั้งค่าการประมวลผล
% Caption: หน้าจอเลือกเครื่องมือและตั้งค่าการประมวลผล
% ==================================================
\begin{figure}[H]
\centering
\includegraphics[width=0.94\textwidth]{\detokenize{images-progress/ภาพ demo หน้าการตั้งค่าการประมวลผล(Tools).jpg}}
\caption{หน้าจอเลือกเครื่องมือและตั้งค่าการประมวลผล}
\label{fig:demo-tools}
\figuredescription{fig:demo-tools}{แสดงหน้าจอเลือก AI Auto-Blur, Detect \& Select หรือ Manual Control พร้อมกำหนดรูปแบบและระดับความเข้มของการปกปิด}
\end{figure}

% ==================================================
% TODO IMAGE: ใส่ภาพ Demo หน้าจอการตรวจจับและเลือกบุคคล
% Caption: หน้าจอการตรวจจับและเลือกบุคคลที่ต้องการปกปิด
% ==================================================
\begin{figure}[H]
\centering
\includegraphics[width=0.94\textwidth]{\detokenize{images-progress/ภาพ demo หน้าเมนู Detect&Select .jpg}}
\caption{หน้าจอการตรวจจับและเลือกบุคคลที่ต้องการปกปิด}
\label{fig:demo-detect-select}
\figuredescription{fig:demo-detect-select}{แสดงผลการตรวจจับเป็น Face Map เพื่อให้ผู้ใช้ตรวจสอบบุคคลที่พบ เลือกบุคคลที่ต้องการปกปิด และกำหนดรูปแบบการปกปิด}
\end{figure}

% ==================================================
% TODO IMAGE: ใส่ภาพ Demo หน้าจอรายการคีย์เฟรมและ Timeline
% Caption: หน้าจอรายการคีย์เฟรมและช่วงเวลาที่ตรวจพบบุคคล
% ==================================================
\begin{figure}[H]
\centering
\includegraphics[width=0.94\textwidth]{\detokenize{images-progress/ภาพ demo หน้าจอแสดงจำนวน Keyframes.jpg}}
\caption{หน้าจอรายการคีย์เฟรมและช่วงเวลาที่ตรวจพบบุคคล}
\label{fig:demo-timeline}
\figuredescription{fig:demo-timeline}{แสดงรายการคีย์เฟรม หมายเลขเฟรม และเวลาที่ตรวจพบบุคคล เพื่อให้ผู้ใช้เลือกตรวจสอบช่วงสำคัญของวิดีโอได้สะดวก}
\end{figure}

% ==================================================
% TODO IMAGE: ใส่ภาพ Demo หน้าจอ Manual Control
% Caption: หน้าจอควบคุมและกำหนดพื้นที่ปกปิดด้วยตนเอง
% ==================================================
\begin{figure}[H]
\centering
\includegraphics[width=0.94\textwidth]{\detokenize{images-progress/ภาพ demo เมนู Manual Control.jpg}}
\caption{หน้าจอควบคุมและกำหนดพื้นที่ปกปิดด้วยตนเอง}
\label{fig:demo-manual-control}
\figuredescription{fig:demo-manual-control}{แสดงหน้าจอ Manual Control สำหรับวาดกรอบครอบบริเวณที่ต้องการปกปิด ปรับขนาดกรอบ และกำหนดให้กรอบติดตามวัตถุในวิดีโอ}
\end{figure}

% ==================================================
% TODO IMAGE: ใส่ภาพ Demo หน้าจอโหมด AI Auto-Blur
% Caption: หน้าจอการทำงานของโหมด AI Auto-Blur
% หมายเหตุ: ไฟล์ "ภาพ demo AI-Auto Blur.jpg" ซ้ำกับภาพ Manual Control
% จึงเว้นตำแหน่งไว้เพื่อป้องกันภาพและคำอธิบายไม่ตรงกัน
% ==================================================
\begin{figure}[H]
\centering
\reportimageplaceholder{[เว้นพื้นที่สำหรับภาพ Demo โหมด AI Auto-Blur]}
\caption{หน้าจอการทำงานของโหมด AI Auto-Blur}
\label{fig:demo-auto-blur}
\figuredescription{fig:demo-auto-blur}{ตำแหน่งนี้ใช้แสดงหน้าจอการทำงานของโหมด AI Auto-Blur ซึ่งตรวจจับและปกปิดใบหน้าโดยอัตโนมัติ เพื่อลดการกำหนดพื้นที่ทีละตำแหน่ง}
\end{figure}

% ==================================================
% TODO IMAGE: ใส่ภาพ Demo หน้าจอตรวจสอบผลลัพธ์ก่อนส่งออก
% Caption: หน้าจอตรวจสอบผลลัพธ์หลังประมวลผลก่อนส่งออกไฟล์
% ==================================================
\begin{figure}[H]
\centering
\includegraphics[width=0.94\textwidth]{\detokenize{images-progress/ภาพ หน้า Preview Result.jpg}}
\caption{หน้าจอตรวจสอบผลลัพธ์หลังประมวลผลก่อนส่งออกไฟล์}
\label{fig:demo-preview}
\figuredescription{fig:demo-preview}{เป็นหน้าจอตรวจสอบผลหลังประมวลผล ผู้ใช้สามารถเล่นวิดีโอตัวอย่างเพื่อยืนยันว่าพื้นที่เป้าหมายถูกปกปิดครบถ้วนก่อนส่งออกไฟล์}
\end{figure}

\section{การพัฒนาและการปรับปรุงจากภาคการศึกษาก่อนหน้า}

ตารางที่ \ref{tab:progress-update} เปรียบเทียบสิ่งที่รายงานเดิมระบุว่าอยู่ระหว่างพัฒนากับสภาพของซอร์สโค้ดปัจจุบัน

\begin{table}[H]
\centering
\caption{ความก้าวหน้าของระบบในรอบปัจจุบัน}
\label{tab:progress-update}
\begin{tabular}{|p{0.29\textwidth}|p{0.25\textwidth}|p{0.31\textwidth}|}
\hline
\textbf{หัวข้อ} & \textbf{สถานะในรายงานเดิม} & \textbf{สถานะปัจจุบัน} \\
\hline
ส่วนติดต่อผู้ใช้ & ออกแบบ UI/UX และทดสอบเชื่อม API & มีกระบวนการอัปโหลด Tools, Face Map, Manual Control, Preview และ Export \\
\hline
Selective Blur & พัฒนาเบื้องต้น & จัดกลุ่มบุคคลจากเวกเตอร์ลักษณะ บันทึกตัวเลือกการปกปิด และเลือกปกปิดหรือยกเว้นรายบุคคลได้ \\
\hline
การประมวลผลวิดีโอ & อยู่ระหว่างพัฒนา & มีกระบวนการตรวจจับแบบเบื้องหลัง การบันทึกตำแหน่งตามเฟรม และบริการสร้างไฟล์ผลลัพธ์ \\
\hline
Manual Override & ออกแบบหน้าจอไว้ & วาด ย้าย ปรับขนาด ลบ และบันทึกกรอบปกปิด พร้อมตัวเลือกติดตามวัตถุ \\
\hline
รูปแบบการปกปิด & กำหนดแนวทางไว้ & ใช้ Gaussian Blur, Pixelation และ Black Box พร้อมปรับระดับความเข้ม \\
\hline
Preview และ Export & ระบุเป็นหน้าผลลัพธ์ & สร้างตัวอย่างก่อนส่งออกและส่งคืน URL สำหรับดาวน์โหลดหรือเปิดผ่าน Cloudinary \\
\hline
การประเมินเชิงปริมาณ & อยู่ระหว่างพัฒนา & ยังต้องเก็บผลการประเมินประสิทธิภาพและความแม่นยำด้วยชุดทดสอบมาตรฐาน \\
\hline
\end{tabular}
\end{table}

\subsection{การตรวจจับและสร้าง Face Map}

เมื่อเริ่มกระบวนการตรวจจับ ระบบอ่านวิดีโอทีละเฟรมและเรียกบริการตรวจจับใบหน้า จากนั้นตัดภาพใบหน้าพร้อมขยายขอบเขตบางส่วนเพื่อสร้างเวกเตอร์ลักษณะ ระบบเปรียบเทียบเวกเตอร์ดังกล่าวกับค่าเฉลี่ยของกลุ่มบุคคลที่พบก่อนหน้า หากค่าความคล้ายสูงกว่าค่าเกณฑ์ ระบบจะรวมใบหน้าเข้ากับกลุ่มเดิม มิฉะนั้นจะสร้างกลุ่มบุคคลใหม่ เมื่อกระบวนการเสร็จสิ้น ระบบจะสร้างรายการ Face Map ค่าความเชื่อมั่นเฉลี่ย และ Timeline ของการปรากฏ

\subsection{การปกปิดแบบเลือกได้และการสร้างผลลัพธ์}

กลุ่มบุคคลที่สร้างขึ้นจะมีสถานะเริ่มต้นเป็นการปกปิด ระบบนำสถานะนี้ไปรวมกับตำแหน่งใบหน้าในทุกเฟรมเมื่อสร้างตัวอย่างหรือส่งออกไฟล์ หากผู้ใช้ต้องการให้ใบหน้าบางรายยังคงมองเห็น สามารถปิดสถานะการปกปิดของบุคคลนั้นได้ บริการสร้างผลลัพธ์จะตรวจสอบขอบเขตพื้นที่ก่อนใช้เอฟเฟกต์ที่ผู้ใช้เลือก และกำหนดขนาด kernel ที่เหมาะสมสำหรับ Gaussian Blur

\subsection{การปกปิดด้วยตนเองและการติดตาม}

Manual Control แสดงภาพตัวอย่างให้ผู้ใช้ลากกรอบเหนือวัตถุที่ต้องการปกปิด ตำแหน่งบนหน้าจอจะถูกแปลงกลับเป็นพิกัดจริงของภาพก่อนบันทึก สำหรับวิดีโอ ระบบจะเริ่มตัวติดตามที่เฟรมเริ่มต้นของกรอบ และนำกรอบที่ติดตามได้ไปใช้ปกปิดเฟรมถัดไป หากตัวติดตามไม่พร้อมใช้งาน ระบบจะใช้กรอบแบบคงที่เป็นทางเลือกสำรอง

\section{การตรวจสอบระบบในรอบจัดทำรายงาน}

ผลการตรวจสอบในตารางที่ \ref{tab:inspection-result} เป็นการตรวจเส้นทางการทำงานและซอร์สโค้ดปัจจุบัน ไม่ใช่ผลการประเมินประสิทธิภาพของแบบจำลองหรือผลการทดสอบกับผู้ใช้จริง จึงแยกสถานะ ``พัฒนาแล้ว'' ออกจาก ``ต้องดำเนินการต่อ'' อย่างชัดเจน

\begin{table}[H]
\centering
\caption{ผลการตรวจสอบระบบในรอบจัดทำรายงาน}
\label{tab:inspection-result}
\begin{tabular}{|p{0.13\textwidth}|p{0.38\textwidth}|p{0.34\textwidth}|}
\hline
\textbf{รายการ} & \textbf{วิธีตรวจสอบ} & \textbf{ผล} \\
\hline
TC-01 Upload & ตรวจสอบ endpoint และชนิดไฟล์ที่ส่วนหลังรองรับ & พัฒนาแล้วสำหรับภาพและวิดีโอที่กำหนด \\
\hline
TC-02 Detection Pipeline & ตรวจสอบลำดับการตรวจจับ การรู้จำ การติดตาม และการบันทึกฐานข้อมูล & พัฒนาแล้วในบริการส่วนหลัง \\
\hline
TC-03 Selective Blur & ตรวจสอบ API และสถานะการเลือกปกปิดรายบุคคล & พัฒนาแล้ว \\
\hline
TC-04 Manual Blur & ตรวจสอบคอมโพเนนต์วาดกรอบและ API สำหรับบันทึกกรอบ & พัฒนาแล้ว พร้อมตัวเลือกการติดตาม \\
\hline
TC-05 Preview/Export & ตรวจสอบ route และบริการสร้างตัวอย่างกับไฟล์ผลลัพธ์ & พัฒนาแล้ว \\
\hline
TC-06 Frontend Lint & รัน \texttt{npm.cmd run lint} & ยังไม่ผ่าน: 37 errors และ 23 warnings โดยส่วนใหญ่เป็น explicit \texttt{any} \\
\hline
TC-07 End-to-end Benchmark & ทดสอบจริงพร้อมแบบจำลองและข้อมูลรับรองของ Cloudinary & ต้องดำเนินการต่อ \\
\hline
\end{tabular}
\end{table}

ผล TC-06 ชี้ว่าควรปรับปรุงความปลอดภัยด้านชนิดข้อมูลและมาตรฐานของ Next.js ก่อนใช้งานจริง ส่วน TC-07 ต้องเตรียมแบบจำลองและข้อมูลรับรอง Cloudinary ให้ครบ เพื่อทดสอบกระบวนการตั้งแต่การอัปโหลดถึงการดาวน์โหลดผลลัพธ์
